% concatenated from jvectoronLCG_-_66.tex
\documentclass{amsart}
\usepackage{amssymb}
\usepackage{amsthm}
\usepackage[dvipsnames]{xcolor}
\usepackage{verbatim}
\usepackage{graphicx}
%\documentstyle[12pt]{article}
% ----------------------------------------------------------------
\vfuzz2pt \hfuzz2pt
% THEOREMS -------------------------------------------------------
\newtheorem{thm}{Theorem}[section]
\newtheorem{cor}[thm]{Corollary}
\newtheorem{lem}[thm]{Lemma}
\newtheorem{prop}[thm]{Proposition}
%\newtheorem{corollary}[thm]{Corollary}
\newtheorem{exam}[thm]{Example}
\theoremstyle{definition}\newtheorem{defn}[thm]{Definition}
\theoremstyle{remark}
\newtheorem{rem}[thm]{Remark}
\numberwithin{equation}{section}
% MATH -----------------------------------------------------------
\newcommand{\norm}[1]{\left\Vert#1\right\Vert}
\newcommand{\abs}[1]{\left\vert#1\right\vert}
\newcommand{\set}[1]{\left\{#1\right\}}
\newcommand{\Real}{\mathbb R}
\newcommand{\eps}{\varepsilon}
\newcommand{\To}{\longrightarrow}
\newcommand{\BX}{\mathbf{B}(X)}
\newcommand{\A}{\mathcal{A}}
\newcommand{\B}{\mathcal{B}}
\newcommand{\R}{\mathcal{R}}
\newcommand{\N}{\mathcal{N}}
\newcommand{\Complex}{\mathbb{C}}
\newcommand{\va}{\varphi}
\begin{document}

\title[$J$-class  weighted translations on L. C. G. ]
{$J$-class weighted translations
 on locally compact groups}

\author{\bf  M. R. Azimi$^{*^1}$,   I.
Akbarbaglu$^2$ and A. R. Imanzadeh Fard$^3$ }


\address {M. R. Azimi}
\email{mhr.azimi@maragheh.ac.ir(for corresponding)}

\address{I. Akbarbaglu}
\email{i.akbarbaglu@cfu.ac.ir}

\address{A. R. Imanzadeh Fard}
\email{a.imanzadeh98@gmail.com}


\address{$1, 3$: Department of Mathematics, Faculty of Sciences,
University of Maragheh, P.O. Box: 55181-83111, Golshahr, Maragheh,
Iran}

\address{ $^{2}$ Department of Mathematics Education,
Farhangian University, P.O. Box: 14665-889, Tehran, Iran}

\thanks{$^*$Corresponding author}



\subjclass[2010]{Primary 47A16; Secondary 47B38.}

\keywords{Extended limit set, $J$-class, Orbit, Hypercyclic,
Weighted translation, Locally compact group, Haar measure.}


\date{}

\dedicatory{}

\commby{}

%%% ----------------------------------------------------------------------

\begin{abstract}
A bounded linear operator $T$ on a Banach space $X$ (not necessarily
separable) is said to be $J$-class operator whenever the extended
limit set, say $J_T(x)$ equals $X$ for some vector $x\in X$.
Practically, the extended limit sets localize the dynamical behavior
of operators.
 In this paper, using the extended limit sets we will examine the necessary
 and sufficient conditions for the  weighted translation $T_{a,\omega}$
  to be $J$-class  on a locally compact group $G$,
 within the setting of $ L^p$-spaces for $ 1 \leq p < \infty $.
 Precisely, we delineate the boundary between $J$-class and
hypercyclic behavior for weighted translations. Then, we will show
that for  torsion elements in  locally compact
 groups,  unlike the case of  non-dense orbits of  weighted
translations, we have $J_{T_{a,\omega}}(0)=L^p(G)$.
 Finally, we will provide some examples on which the  weighted translation
$ T_{a,\omega}$ is $J$-class but it fails to be hypercyclic.



\end{abstract}
%%% ----------------------------------------------------------------------
\maketitle
%%% ---
\section{\textbf{Introduction}}
Let's assume that $ G $ is a locally compact group with an identity
element $ e $ and equipped with the right Haar measure $\lambda$.
 An element $ g \in G $ is called a torsion element if there exists
 $ n \in \mathbb{N} $ such that $ g^n = e $ and called
 periodic if the closed subgroup generated by $g$, is compact. It is clear
 that every element of torsion (finite order) is also periodic, and we
 call that $ g $ aperiodic if it is not periodic.
We say that an element $a\in G$ passes through  compact subsets,
 whenever for each compact subset $K\subset G$,  there exists an $N \in
\mathbb{N}$ such that $K \cap Ka^{\pm m}= \emptyset$, for each  $ m
\geq N $. Apparently, by \cite[Lemma 2.1]{chu},   every  aperiodic
element of the second countable groups has this property.

 For an arbitrary element
 $g \in G $, the unit point mass  measure (or Dirac measure) at $g$, denoted $\delta_g$,   for any
  subset $A \subset G $ is defined by  $ \delta_g(A) = 1 $ if
 $ g \in A $ and $ \delta_g(A) = 0$, otherwise.
 The Banach space $ L^p(G) $ (for $ 1 \leq p < \infty $) consists
of all Borel functions $ f: G \rightarrow \mathbb{C}$ such that:
\[
\|f\|_p = \left( \int_G |f(x)|^p \, d\lambda(x)
\right)^{\frac{1}{p}} < \infty.
\]
Let a continuous function $ \omega: G \to (0, \infty) $ be a
\emph{weight} on a group $ G $. Given an element $ a \in G $, we
define  typically the \emph{weighted translation operator} $
T_{a,\omega}: L^p(G) \to L^p(G) $ for $ 1 \leq p < \infty $ by the
assignment:
 $$ f\mapsto \omega f\ast\delta_a,\qquad\forall f\in L^p(G),$$

where, for every $ x \in G$, the convolution:
\begin{eqnarray*}
(f\ast\delta_a)(x)=\int_G\,\,f(xy^{-1})\,d\delta_a(y)=f(xa^{-1})
\end{eqnarray*}
is interpreted as  the right translation of $f$ by $a^{-1}$. For
each Borel-measurable subset $A$, the \emph{change of variable
formula},  $$\int_{A} f*\delta_a(x)d\lambda(x)=
  \int_{Aa^{-1}} f(x) d\lambda(x),$$ will be used
  frequently throughout this paper. For
each $m\in \mathbb{N}$, the iterates of the weighted translation
 $T_{a,\omega}$ is obtained as follows:
\begin{eqnarray*}
T^m_{a,\omega}(f)(x)=[\prod_{i=0}^{i=m-1}\omega(xa^{-i})]f(xa^{-m})
\end{eqnarray*}
For simplicity of the notation, we adopt that:
 $${\tilde{\omega}_{m}}(x)=(\prod_{i=0}^{i=m-1}\omega(xa^{-i}))^{-1},
 \qquad \omega_{m}(x)=\prod_{i=1}^{i=m}\omega(xa^{i}).$$
 So
\begin{eqnarray*}
T^m_{a,\omega}(f)(x)={\tilde{\omega}_{m}}^{-1}(x)f(xa^{-m}).
\end{eqnarray*}


Let $X$ be a separable Banach space. A bounded linear operator $T:
X\rightarrow X$ is said to be hypercyclic whenever there exists a
vector $x\in X$ such that its orbit, $$orb(T,x)=\{T^nx:
n=0,1,2,\cdot\cdot\cdot\}$$ is dense in $X$. Concerning dynamics of
linear operators, the interested readers are referred to the books
\cite{baya, gros}.  In the year 1995, H. Salas was the first
mathematician who has initiated the hypercyclic weighted shifts in
terms of their weight sequences in \cite{sal}. In \cite{chen},  C.
Chen and C-H. Chu have studied the hypercyclicity of the weighted
translation $T_{a, \omega}$ on locally compact groups for various
cases of the element $a$. They provided a necessary and sufficient
conditions for the hypercyclicity of $T_{a, \omega}$ when $a$ is an
aperiodic element. Afterwards,   G. Costakis and  A. Manoussos have
localized the concept of hypercyclicity in \cite{cos}. For this,
they introduced the limit sets and the extended limit sets. Given a
bounded linear operator $T$ on a Banach space $X$ (not necessarily
separable) and $x\in X$, the \emph{limit set} is defined as follows:
$$L_T(x)=\{y\in X | ~ \exists (n_k) \subseteq \mathbb{N},~
\text{s.t.} T^{n_k} x \to y \}.$$
 Note that $ T $ is hypercyclic if there exists a vector
$ x \in X $ such that $L_T(x)=X$.  Additionally, we denote the
\emph{extended limit set} by $ J_T(x) $. It
 consists of all vectors $y$ in $X$ on which there exist a strictly increasing
  sequence of positive integers ${n_k}$ and a sequence $(x_n)\subseteq X$ such that
   $x_n \to x$  and $T^{n_k} x_n \to y$.  Shortly,
  \[
   J_T(x) = \{ y \in X :\exists \upuparrows(n_k) \subseteq \mathbb{N},   \exists (x_n)
   \subset X ~ \text{s.t. }~ ~ x_n \to x \text{ and } T^{n_k} x_n \to y \}
   .\]
An operator $ T $ is said to be  \textit{J-class } whenever $ J_T(x)
= X $ for some vector $x\in X$. In this case, a vector $ x $ is
called a \textit{J-vector} for $ T $. For all $x\in X$, the sets $
L_T(x)$  $ J_T(x)$ are  closed and $T$-invariant. In the case, an
operator $T$ is power bounded i.e., for each $n\in \mathbb{N}$,
$\|T^n\|\leq M$ (for some real number) then $ J_T(x)=L_T(x)$
(\cite{man}).


 It is worth noting that if an
operator $T$ is hypercyclic then $ J_T(x) = X $ for each vector
$x\in X$. Further,  an equivalent definition of $J_{T}(x)$ can be
stated by the way that $J_{T}(x)$ consists of those $y\in X$ such
that for every pair of neighborhoods $ U, V $ of $ x, y$
respectively and for each positive integer $N$, there exists a
positive integer $ n> N$ such that $$T^n U \cap V\neq \emptyset.$$
Note that on a separable Banach space $X$, an operator $T$ is
topologically transitive if and only if $J_{T}(x)=X$ for some cyclic
vector $x\in X$, see \cite{man} for more details. In \cite{cos}, G.
Costakis and A. Manoussos have characterized the $J$-class
($J^{mix}$- class) unilateral weighted shifts on
$\ell^{\infty}(\mathbb{N})$ in terms of their weight sequences
through the use of the extended limit sets. Also they proved that a
bilateral weighted shift on $\ell^{\infty}(\mathbb{Z})$ cannot be a
$J$-class. The existence of the $J$-class operators on Banach spaces
has been studied in \cite{nas2}. Some other works like as \cite{as,
az3, azmu, mar} are devoted to the study on $J$-class operators as
well.

Throughout this paper, the interior and the closure of any subset
$A$ will be denoted by $A^\circ$ and $\overline{A}$, respectively.
By $\sigma(f)$ and $ess\,sup f$ we mean the support of any function
$f\in L^p(G)$ i.e., $\{x\in G: \, f(x)\neq 0\}$ and its essential
supremum norm, respectively. The set of all continuous functions on
$G$ with compact supports will be denoted by $C_c(G)$ as well.

In this paper, first we will show that for the torsion elements of
the locally compact
 groups,  weighted translations  cannot be  $J$-class operators.
Then by imposing special conditions on the weight function, we aim
to examine the $J$-class weighted translations $ T_{a,\omega}$ on
$L^p(G)$. Moreover, we will provide some examples on which
$T_{a,\omega}$ is $J$-class but it fails to be hypercyclic.





\section{\textbf{$J$-class  weighted translation operators on  $ L^p(G) $}}
 %%%%%%%%%%%%%%%%%%%%%%%%%%%%%%%%%%%%%%%%%%%%
 In this section,  we shall provide the necessary and sufficient conditions for
 weighted translation operators to  possess nonzero $J$-vectors
  in the case of non torsion elements. Furthermore,
 we show that for the torsion elements of  locally compact
 groups, weighted translation operators unlike hypercyclicity
 (\cite[Lemma 1.1]{chu}), the extended limit set  $J_{T_{a,\omega}}(0)$
  can fill whole the space $L^p(G)$.


%%%%%%%%%%%%%%



\begin{thm}\label{t2.2}
Let $G$ be a locally compact group, $a\in G$ be an element passing
through  each compact subsets and $\omega:G\to (0,\infty)$ be a
weight function. If the operator $T_{a,\omega}$   is $J$-class then
for each compact subset $\Delta \subset G$ with the positive
measure, there is a sequence of subsets $(E_k)$ such that
$\lambda(E_k)\to\lambda(\Delta)$. Moreover, for some subsequence
$(n_k)$,
$$
\lim_{k\rightarrow \infty}\underset{x\in E_{k}}{ess\,sup}\,
 \tilde{\omega}_{n_k}(x)=0.
$$
\end{thm}
\begin{proof}
Suppose that $T_{a,\omega}$ is a  $J$-class operator. Then there
exists a nonzero $f \in L^p(G)$ such that
$J_{T_{a,\omega}}(f)=L^p(G)$. Let $0<\epsilon<1$ and let $\Delta$ be
a compact subset with positive measure. Choose $0<\eta<1$ with
$4\eta^p<\lambda(\Delta)$ and $\frac{2\eta}{1-\eta}<\epsilon$ so
that the set $A=\{x\in\sigma(f):|f(x)|>\eta\}$ is of positive finite
measure. By regularity of Haar measure $\lambda$ we may find a
compact subset $B\subseteq A$ such that $\lambda(A\setminus
B)<\eta^p$. Also, pick $0<\delta<\eta$ and $n\in\mathbb{N}$ such
that $(\Delta\cup B)a^{-n}\cap(\Delta\cup B)=\emptyset$ and moreover
$\|f-g\|_p<\delta^2$ and
$\|T_{a,\omega}^n(g)-\chi_{\Delta}\|_p<\delta^2$. Define $C=\{x\in
G\setminus B: |f(x)-g(x)|\geq\delta\}$. Now either $\lambda(C)=0$ or
$\lambda(C)<\delta^p$, because of
$$
\delta^{2p}>\|f-g\|_{p}^p\geq\int_C|f(x)-g(x)|^p
d\lambda(x)\geq\delta^p\lambda(C).
$$
Also, set
$D=\{x\in\Delta:|\tilde{\omega}_n(x)^{-1}g(xa^{-n})-1|\geq\delta\}$.
In the case $D\neq \emptyset$, we have
$$
\delta^{2p}>\|T_{a,\omega}^n(g)-\chi_{\Delta}\|_{p}^p\geq\int_D|\tilde{\omega}_n(x)^{-1}g(xa^{-n})-1|^p
d\lambda(x)\geq\delta^p\lambda(D).
$$
Now, for $x\in E:=\Delta\setminus \big(D\cup Ca^n\cup(A\setminus
B)a^n\big)$ we have $\lambda(\Delta\setminus E)<3\eta^p$ and
furthermore
$$
\tilde{\omega}_n(x)<\frac{|g(xa^{-n})|}{1-\delta}
<\frac{|f(xa^{-n})|+\delta}{1-\delta}<\frac{2\eta}{1-\eta}<\epsilon.
$$
\end{proof}
\begin{rem}
It is worth noting that on the support of the nonzero $J$-vector
$f$, set $f=\lim_{m\to\infty}f\chi_{A_m}$ where $A_m=\{x\in
\sigma(f): |f(x)|\geq\frac{1}{m}\}$. Pick $g\in L^p(G)$ and
$n\in\mathbb{N}$ on which $\|f-g\|_p<\delta^2$ and
$\|T_{a,\omega}^{n}g\|<\delta^2$. Set $B_m=\{x\in A_m:
|f(x)-g(x)|\geq\delta\}$. Then similarly
$0\leq\lambda(B_m)<\delta^p$. Also, by implementing the change of
variable formula to the subset   $C_m=\{x\in A_m:
|\omega_n(x)g(x)|\geq\delta\}$ we get that
$0\leq\lambda(C_m)<\delta^p$. Finally, for every $x\in E_m:=
A_m\setminus(B_m\cup C_m)$ these are inferred that
$\lambda(A_m\setminus E_m)<2\delta^p$ and
$$
\omega_n(x)<\frac{\delta}{|g(x)|}
<\frac{\delta}{\frac{1}{m}-\delta}<\epsilon.
$$
\end{rem}
In the following theorem sufficient conditions for $J$-class
weighted translations on $L^p(G)$-spaces are characterized.
Moreover, it precisely delineates the boundary between $J$-class and
hypercyclic behavior for such operators.
\begin{thm}\label{t4}
Let $G$ be a locally compact group, $a\in G$ be an element passing
through  each compact subsets and $\omega:G\to (0,\infty)$ be a
weight function. If
\begin{itemize}
  \item [(i)]for each compact subset $\Delta \subset G$ with
the positive measure,  there is a sequence of Borel subsets $(E_k)$
such that $\lambda(E_k)\to\lambda(\Delta)$. Moreover, for some
subsequence $(n_k)$, $\lim_{k\rightarrow \infty}\underset{x\in
E_{k}}{ess\,sup}\,
 \tilde{\omega}_{n_k}(x)=0$;
  \item [(ii)] for some compact subset $K \subset G$ with
$\lambda(K)>0$,
 $$ \lim_{k\rightarrow \infty}\underset{x\in
K}{ess\,sup}\,
 {\omega}_{n_k}(x)=0.$$
\end{itemize}
Then the operator $T_{a,\omega}$ is $J$-class with the $J$-vector
$\chi_K$.
\end{thm}
\begin{proof}
Assume that the statements of $\mbox{(i)}$ and $\mbox{(ii)}$ hold. We shall prove
that $J_{T_{a,\omega}}(\chi_K)=L^p(G)$. Note that
$\overline{C_c(G)}=L^p(G)$.
 For an arbitrary $f\in C_c(G)$ it is known that
$\sigma(f)$ is a non-empty and compact subset of $G$. We are going to find a
sequence like as $(g_k)$ in $L^p(G)$ and a strictly increasing sequence of positive integers $(n_k)$ such that
$\|g_k -\chi_K\|_p<\epsilon$ and $\|T_{a,\omega}^{n_k}(g_k)-f\|_p<\epsilon$.
Now suppose that $E_k\subseteq \sigma(f)$, with $\lim_{k\rightarrow
\infty}\underset{x\in E_{k}}{ess\,sup}\,
 \tilde{\omega}_{n_k}(x)=0$.  Since $a$ passes through  each compact subset,
  there exists an $m\in\mathbb{N}$ such that $\sigma(f) \cap
\sigma(f)a^{\pm n}=\emptyset$ for all $n>m$.
 Moreover, for each $\epsilon>0$, our hypotheses  imply that
$\lambda(\sigma(f)-E_k)<\frac{\epsilon} {2^p \Vert
f\Vert_{\infty}^{p}}$,
$ess\,sup\,\tilde{\omega}_{n_k}<\frac{\epsilon}{\|f\|_{p}^p}$ on $E_k$
and $ess\,sup\,\omega_{n_k}<\frac{\epsilon}{2^p\|f\|_{p}^p}$ on $K$
for sufficiently large $k$ and $n_k$. Now, define
$$g_k(x):=\tilde{\omega}_{n_k}(xa^{n_k})f\chi_{E_k}(xa^{n_k})+\chi_{K}(x).$$ Then, by
the change of variable formula and the aforementioned facts we have
\begin{eqnarray*}
 \|g_k-\chi_{K}\|_p^p&=&\int_G|\tilde{\omega}_{n_k}(xa^{n_k})
 f\chi_{E_k}(xa^{n_k})|^p d\lambda(x) \\
  &=&\int_G|\tilde{\omega}_{n_k}(x)f\chi_{E_k}(x)|^p d\lambda(x) \\
&=&\int_{E_k}|\tilde{\omega}_{n_k}(x)f(x)|^p d\lambda(x)\\
      &<& \epsilon.
  \end{eqnarray*}
Therefore, $\Vert g_{k}-\chi_{K}\Vert_p\to 0$ as $k\to \infty$. On
the other hand, we have
\begin{eqnarray*}
&&\Vert T_{a,\omega}^{n_k}(g_k) - f\Vert_p^p
\\&=&
\int_G \vert T_{a,\omega}^{n_k}(g_k)(x) -f(x)\vert^p d\lambda(x)
\\&=&
\int_G \vert
\omega(x)\omega(xa^{-1})\cdot\cdot\cdot\omega(xa^{-n_k+1})g_k(xa^{-n_k})
-f(x)\vert^p d\lambda(x)
\\&=&
\int_G \vert
\omega(x)\omega(xa^{-1})\cdot\cdot\cdot\omega(xa^{-n_k+1})
\tilde{\omega}_{n_k}(x)f(x)\chi_{E_k}(x)\\&& +
\omega(x)\omega(xa^{-1})\cdot\cdot\cdot\omega(a^{-n_k+1}x)\chi_K(xa^{-n_k})
-f(x)\vert^p d\lambda(x)
\\&\leq&
2^p\int_G|f\chi_{E_k}(x)-f(x)|^p d\lambda(x)
\\&&+
2^p\int_G
\vert\omega(x)\omega(a^{-1}x)\cdot\cdot\cdot\omega(xa^{-n_k+1})\chi_K(xa^{-n_k})
\vert^p d\lambda(x)
\\&\leq&
2^p\lambda\big(\sigma(f)-E_k\big)\|f\|_{\infty}^p
+
2^p\int_G \vert{\omega}_{n_k}(x)\chi_K(x) \vert^p d\lambda(x)\\&<& \epsilon.
\end{eqnarray*}
\end{proof}
\begin{thm}
Let $G$ be a locally compact group with the  right Haar measure
$\lambda$ and let $a\in G$ be a torsion element of order $\gamma$.
Then the following conditions are equivalent:
\begin{itemize}
\item[(i)] $J_{T_{a,\omega}}(0)=L^p(G)$;\\
\item[(ii)] For every compact set $F\subseteq G$ of positive measure and
every $\epsilon, \delta>0$ there exist $n\in\mathbb{N}$ and a Borel
subset $E\subseteq F$ such that $\lambda(F\setminus E)<\delta$ and
$$
\underset{x\in E}{ess\,sup}\, \omega_n(x)^{-1}<\epsilon.
$$
\end{itemize}
\end{thm}
\begin{proof}
To prove the implication $\mbox{(i)}\Rightarrow\mbox{(ii)}$, suppose
that  $J_{T_{a,\omega}}(0)=L^p(G)$. Take $\epsilon, \delta>0$ and a
compact subset $F\subseteq G$ arbitrarily. Without loss of
generality, assume that $\lambda(F)>3\delta^{p}$ and
$\frac{\delta}{1-\delta}<\epsilon$. Hence, we may find $g\in L^p(G)$
and $n\in\mathbb{N}$ large enough satisfying $\|g\|_p<\delta^2$ and
$$
\|T^n_{a,\omega}g-\chi_{\bigcup_{k=0}^{\gamma-1}Fa^k}\|_p<\delta^2.
$$
Set $A=\{x\in F: |g(x)|\geq\delta\}$. In the case $A$ is a non-empty
subset, it is inferred that
$$
\delta^{2p}>\|g\|_{p}^p\geq\int_A |g(x)|^p
d\lambda(x)\geq\delta^p\lambda(A),
$$
that is $\lambda(A)<\delta^p$. Also, set
$$
B=\left\{x\in F:\,\left|\omega_n(x)g(x)-1\right|\geq\delta\right\}.
$$
Now, we may write $n=\gamma[\frac{n}{\gamma}]+k$ where $0\leq
k\leq\gamma-1$. If $\lambda(B)\neq 0$ then by the  change of
variable formula  we obtain that
\begin{eqnarray*}
\delta^{2p} &>&
\|T^{n}_{a,\omega}g-\chi_{\bigcup_{k=0}^{\gamma-1}Fa^k}\|_{p}^p\\
&=& \int_G |\omega(x)\omega(xa^{-1})\cdots
\omega(xa^{-(n-1)})g(xa^{-n})-\chi_{\bigcup_{k=0}^{\gamma-1}Fa^k}(x)|^p d\lambda(x)\\
&=& \int_G |\omega(xa^{n})\omega(xa^{n-1})\cdots
\omega(x a)g(x)-\chi_{\bigcup_{k=0}^{\gamma-1}Fa^k}(xa^n)|^p d\lambda(x)\\
&\geq&
\int_B |\omega_n(x)g(x)-1|^p d\lambda(x)\\
&\geq& \delta^p\lambda(B).
\end{eqnarray*}
Therefore, $\lambda(B)\leq\delta^p$ and subsequently
$\lambda\big(F\setminus(A\cup B)\big)>\delta^p$. Finally, by setting
$E:=F\setminus(A\cup B)$ we obtain $\lambda(F\setminus E)<2\delta^p$
and in turn for each $x\in E$ we have
$$
\frac{1}{\omega_n(x)}<\frac{|g(x)|}
{1-\delta}<\frac{\delta}{1-\delta}<\varepsilon.
$$


To  prove the converse implication
$\mbox{(ii)}\Rightarrow\mbox{(i)}$, the initial definition of
$J_{T_{a,\omega}}(0)$ does not use here any longer. Define the
operator $S_{a,\omega}:L^p(G)\rightarrow L^p(G)$ by
$S_{a,\omega}(h)(x)=\frac{h(xa)}{\omega(xa)}$ for each $x\in G$ and
$h\in L^p(G)$. Take open neighborhood $U$ of the origin and a
non-empty open neighborhood $V$ in $L^p(G)$  arbitrarily. As the set
$C_c(G)$ is dense in $L^p(G)$,  there is a $g\in C_c(G)\cap V$ and
hereon let $F=\sigma(g)$. Since $\omega$ is continuous, then for
some $\eta>0$, $\inf\{\prod_{i=1}^{k}\omega(xa^{i})^{-1}: 1\leq
k\leq\gamma-1\}>\eta$ for all $x\in F$. Now by our assumptions in
$\mbox{(ii)}$, for each $\epsilon,\delta>0$, there exist $E\subseteq
F$ and a positive integer $n$ satisfying $\lambda(F\setminus
E)<\delta$ and $\underset{x\in
E}{ess\,sup}\,\big(\prod_{i=1}^{n}\omega(xa^{i})\big)^{-1}<\epsilon$.
In particular, if we consider
$\omega_{\gamma}(x)=\omega(x)\omega(xa)\cdots\omega(xa^{\gamma-1})$,
then apparently $\underset{E}{ess\,sup}\,\omega_{\gamma}>1$. We may
also write  $n=\gamma[\frac{n}{\gamma}]+k$ where $0\leq
k\leq\gamma-1$, which implies  that
\begin{eqnarray*}
\|S_{a,\omega}^{\gamma n}(g\chi_{E})\|_{p}^{p} &=&
\int_G\left|\frac{g(xa^{\gamma n})\chi_{E}(xa^{\gamma n})}{\omega(x a)\cdots\omega(xa^{\gamma n})}\right|^p d\lambda(x)\\
&=&
\int_{E}\left|\frac{g(x)}{\omega(x a)\cdots\omega(xa^{\gamma n})}\right|^p d\lambda(x)\\
&\leq&
\underset{}{ess\,sup}\,\omega_{\gamma}^{-n}\,\underset{}{ess\,sup}\,
g\,\lambda(E).
\end{eqnarray*}
Moreover,
$$
\|T_{a,\omega}^{\gamma n}\big(S_{a,\omega}^{\gamma
n}(g\chi_{E})\big)-g\|_p= \|g\chi_{E}-g\|_p\leq
ess\,sup\,g\lambda(F\setminus E)<\delta \,{ess\,sup}\,g.
$$
The last ensures  that $T_{a,\omega}^{\gamma n}(U)\cap
V\neq\emptyset$. On the other hand, in the case
${ess\,sup}\,\omega_{\gamma} \leq 1$, it should be noted that
$$
\|T^{n}_{a,\omega}g\|_{p}^{p}\leq
ess\,sup\,\omega_{\gamma}^{p[\frac{n}{\gamma}]}\int_G
\omega^p(ax)\cdots\omega^p(a^{k}x)|g(x)|^p d\lambda(x)\leq
ess\,sup\,\omega_{\gamma}^{p([\frac{n}{\gamma}]+k)}\|g\|_p^p.
$$
Hence, $\{T^{n}_{a,\omega}g\}_{n=1}^{\infty}$ is bounded for every
$n\in\mathbb{N}$ and $g\in L^p(G)$. Actually, in this case,
$T_{a,\omega}$ is power bounded and then by \cite[Proposition
2.10]{man}, $J_{T_{a,\omega}}(0)=L_{T_{a,\omega}}(0)\neq L^p(G)$.
\end{proof}
\begin{cor}
Let $a$ be a torsion element of a locally compact group $G$. Then
$0\in J_{T_{a,\omega}}(0)^\circ$ if and only if
$J_{T_{a,\omega}}(0)=L^p(G)$.
\end{cor}
\begin{proof}
Suppose that  $0\in J_{T_{a,\omega}}(0)^\circ$. Then, there is a
positive real number $r$ such that $B_{r}(0)\in
J_{T_{a,\omega}}(0)$. Hence we can choose $\delta>0$ sufficiently
small with $2\delta\chi_{\cup_{k=0}^{\gamma-1}Fa^k}\in B_{r}(0)$. It
follows that there exist $g\in L^p(G)$ and $n\in\mathbb{N}$ such
that $\|g\|_p<\delta^3$ and
$\|T^n_{a,\omega}g-2\delta\chi_{\cup_{k=0}^{\gamma-1}Fa^k}\|_p<\delta^2$.
Therefore, similar to the proof of the preceding theorem the rest is
done.
\end{proof}
\begin{thm}\label{t3}
Let $G$ be a  locally compact group. Assume that  an element $a\in
G$ passes through  each compact subsets.  Then the followings are
equivalent:
\begin{itemize}
  \item [(i)] $J_{T_{a,\omega}}(0)=L^p(G)$,
  \item [(ii)] for each compact subset $\Delta
\subset G$ with the positive measure,
  there is a sequence of Borel subsets $(E_k)$ such
that $\lambda(E_k)\to\lambda(\Delta)$ and moreover for some
subsequence $(n_k)$,
$$
\lim_{k\rightarrow \infty}\underset{x\in E_{k}}{ess\,sup}\,
\tilde{\omega}_{n_k}(x)=0.
$$
\end{itemize}
\end{thm}
\begin{proof}
$\mbox{(i)}\Rightarrow \mbox{(ii)}$. Suppose that $J_{T_{a,\omega}}(0)=L^p(G)$.
Hence, for each $\epsilon>0$ and a compact subset $\Delta\subset G$ with positive measure,
we may find a sequence $(g_k)$ in $L^p(G)$ and a strictly increasing sequence of
positive integers $(n_k)$ such that $$\|g_k \|_p <\epsilon ~ \text{
and}~  \|T_{a,\omega}^{n_k}(g_k)-\chi_\Delta\|_p<\epsilon.$$
By the assumption
there exists an $N\in\mathbb{N}$ such that $\Delta \cap \Delta
a^{\pm n}=\emptyset$ for all $n>N$. Without loss of generality, we
may assume that $n_k>N$ for each $k\in \mathbb{N}$.
 Fix  $ \eta> 0$, such that $\frac{\eta}{1-\eta}<\epsilon$
and set $B_{\eta,k}=\{x\in G: |g_k(x)|\geq\eta\}$.
Hence, $0\leq\lambda(B_{\eta,k})<\eta^p$. On the other hand, assume that
$$C_{\eta,k}=\{x\in\Delta:|\tilde{\omega}_{n_k}(x)^{-1}g(xa^{-n_k})-1|\geq\eta\}.$$
In the case $C_{\eta,k}\neq \emptyset$, by the
change of variable formula  we obtain $0\leq\lambda(C_{\eta,k})<\eta^p$.
Finally by setting $E_k:=\Delta\setminus \big(C_{\eta,k}\cup
B_{\eta,k}a^{n_k}\big)$, the desired result is established. Indeed,
we have $\lambda(\Delta\setminus E_k)<\eta^p$ and
$$
\tilde{\omega}_{n_k}(x)<\frac{|g(xa^{-n_k})|}{1-\eta}
<\frac{\eta}{1-\eta}<\epsilon.
$$




$\mbox{(ii)}\Rightarrow \mbox{(i)}$. We shall prove
that $J_{T_{a,\omega}}(0)=L^p(G)$. Indeed, for each $\epsilon>0$ and
$f\in L^p(G)$, we are going to find a sequence
 $(g_k)$ in $L^p(G)$ and a strictly increasing sequence of positive integers
  $(n_k)$  such that $\|g_k
\|_p<\epsilon$ and
$\|T_{a,\omega}^{n_k}(g_k)-f\|_p<\epsilon$.
 Note that
$\overline{C_c(G)}=L^p(G)$ and so  we may prove the last assertion
for each $f \in C_c(G)$.
 It is known that
$\sigma(f)$ is a non-empty and compact subset, hence by the
assumptions, for each $\epsilon>0$,  there exist a Borel subset
$E_k\subseteq \sigma(f)$  such that
$\lambda(\sigma(f)-E_k)<\frac{\epsilon} { \Vert f\Vert_{\infty}^{p}}$ and
 $\underset{x\in E_{k}}{ess\,sup}\,
\tilde{\omega}_{n_k}(x)<\frac{\epsilon}{\Vert f\Vert_p^p}$,
 for enough large positive integers $k$ and $n_k$.
Define
$$g_k(x):=\tilde{\omega}_{n_k}(xa^{n_k})f\chi_{E_k}(xa^{n_k}).$$ Then,   by the
change of variable formula and the aforementioned facts we obtain
\begin{eqnarray*}
 \|g_k\|_p^p&=&\int_G|\tilde{\omega}_{n_k}(xa^{n_k})
 f\chi_{E_k}(xa^{n_k})|^p d\lambda(x) \\
  &=&\int_G|\tilde{\omega}_{n_k}(x)f(x)\chi_{E_k}(x)|^p d\lambda(x) \\
%&=&\int_{\sigma(f)-E_k}|\tilde{\omega}_{n_k}(x)f(x)|^p d\lambda(x)\\
&=& \int_{E_k}|\tilde{\omega}_{n_k}(x)f(x)|^p d\lambda(x)\\
      &<& \epsilon.
  \end{eqnarray*}
Therefore, $\Vert g_{k}\Vert_p\to 0$ as $k\to \infty$. On the other
hand, we have
\begin{eqnarray*}
&&\Vert T_{a,\omega}^{n_k}(g_k) - f\Vert_p^p
\\&=&
\int_G \vert T_{a,\omega}^{n_k}(g_k)(x) -f(x)\vert^p d\lambda(x)
\\&=&
\int_{G} \vert
\omega(x)\omega(xa^{-1})\cdot\cdot\cdot\omega(xa^{-n_k+1})\tilde{\omega}_{n_k}(x)
f\chi_{E_k}(x)-f(x)\vert^p d\lambda(x)\\
&=&\int_{G} \vert f\chi_{E_k}(x)-f(x)\vert^p d\lambda(x)\\ &\leq&
\lambda\big(\sigma(f)-E_k\big) \|f\|_{\infty}^p\\&<&\epsilon.
\end{eqnarray*}
\end{proof}
%%%%%%%%%%%%%%%%%%%%%%%%%
\begin{rem}
Note that in the implication $(ii)\Rightarrow (i)$ of Theorem
\ref{t3}, apparently it has not been assumed that a group  $G$ is to
be second countable. As a result, $J_{T_{a,\omega}}(0)$ can fill the
non-separable $L^p(G)$ spaces while it is impossible for the
hypercyclic phenomena.
%Moreover, this theorem generalizes Corollary
%3.4 obtained in \cite{cos}.
\end{rem}
\begin{thm}
Let $G$ be a locally compact group with a fixed right Haar measure
$\lambda$. Assume that $a\in G$ passes through each compact subsets
of $G$. Then the followings are equivalent:
\begin{itemize}
\item[(i)] $0\in J_{T_{a,\omega}}(0)^\circ$;\\
\item[(ii)] $J_{T_{a,\omega}}(0)=L^p(G)$.
\end{itemize}
\end{thm}
\begin{proof}
It is adequate to prove $\mbox{(i)}\Rightarrow\mbox{(ii)}$, as the
converse is obvious. Take $r>0$ such that the open ball $B_{r}(0)$
is contained in $J_{T_{a,\omega}}(0)$. Pick an arbitrary compact
subset $K\subset G$ with $\lambda(K)>0$ and choose $\epsilon>0$ with
$\epsilon<r$ and $\lambda(K)>3\epsilon^p$. Hence
$\displaystyle\frac{\epsilon}{2}\lambda(K)^{-\frac{1}{p}}\chi_{K}\in
B_{\epsilon}(0)$. It follows that there exist $g\in L^p(G)$ and
$n\in\mathbb{N}$ large enough, so that $K\cap Ka^{-n}=\emptyset$,
$$
\|g\|_p<\frac{\epsilon^3}{4\lambda(K)^\frac{1}{p}} \qquad \mbox{and}
\qquad
\|T_{a,\omega}^n(g)-\frac{\epsilon}{2\lambda(K)^\frac{1}{p}}\chi_{K}\|_p<
\frac{\epsilon^2}{4\lambda(K)^\frac{1}{p}}.
$$
By defining $C=\{x\in G\setminus
K:\,|g(x)|>\frac{\epsilon^2}{4\lambda(K)^\frac{1}{p}}\}$ we get that
$$
\frac{\epsilon^{3p}}{4^p\lambda(K)}>\int_G|g(x)|^p\,d\lambda(x)\geq
\frac{\epsilon^{2p}}{4^p\lambda(K)}\lambda(C).
$$
Therefore, in the case $\lambda(C)\neq 0$, it is inferred that
$0<\lambda(C)<\epsilon^p$. Also, put
$$
D=\left\{x\in K:\left|\prod_{i=0}^{n-1}
\omega(xa^{-i})g(xa^{-n})-\frac{\epsilon}{2\lambda(K)^{\frac{1}{p}}}\right|>
\frac{\epsilon}{4\lambda(K)^{\frac{1}{p}}}\right\}.
$$
It implies that if $\lambda(D)\neq 0$, then
$$
\frac{\epsilon^{2p}}{4^p\lambda(K)}>
\int_G|T_{a,\omega}^{n}g(x)-\frac{\epsilon}{2\lambda(K)^\frac{1}{p}}\chi_{K}(x)|^p
d\lambda(x)\geq \frac{\epsilon^{p}}{4^p\lambda(K)}\lambda(D),
$$
whence $0<\lambda(D)<\epsilon^p$. Now for each $x\in
K\setminus(D\cup Ca^{n})$ we obtain that
$$
\frac{1}{\prod_{i=0}^{n-1}\omega(xa^{-i})}<\frac{4\lambda(K)^{\frac{1}{p}}}
{\epsilon}|g(xa^{-n})|<\epsilon.
$$
\end{proof}


%%%%%%%%%%%%%%%%%%%%%%%%%%%%%%%%%%%%%%%%%%%
\begin{exam}
Let $G=(\mathbb{R}, +)$, be the group of the real numbers equipped
with  Lebesgue measure.  Define the weight function $w$ on
$\mathbb{R}$ by
$$w(x)=\left\{
  \begin{array}{ll}
   \alpha, & {1\leq x}, \\
        -\frac{x}{2}+1, & {-1< x < 1},\\
    \beta, &{x\leq -1} ,
  \end{array}
\right.$$  where the real numbers $\alpha,\beta$ are arbitrarily
chosen such that $1<\alpha <\beta$. Fix a  nonzero number $a\in
\mathbb{R}$ and consider the convolution
$(\omega*\delta_{a})(x)=\omega(x- a)$, for each $x\in \mathbb{R}$
and $f\in L^p(\mathbb{R})$.  The statement of Theorem \ref{t3}
ensures that $J_{T_{a,\omega}}(0)=L^p(\mathbb{R})$ while is not
hypercyclic by \cite[Theorem 2.2]{chu}.
 \end{exam}
\begin{exam}
Let's assume that $ G = (\mathbb{R}^+,\times) $ and $ a =\frac{1}{2}
$ with the weight function $ \omega(x) = \exp(x)$. We will show that
$J_{T_{a,\omega}}(0)=L^p(\mathbb{R})$. For this, consider that
 \begin{eqnarray*}
\prod_{i=0}^{i=m-1}\omega(xa^{-i})&=&\prod_{i=0}^{i=m-1}\omega(2^{i}x)
 =\prod_{i=0}^{i=m-1}\exp(2^{i}x)=\prod_{i=0}^{i=m-1}\exp(2^{i}x) \\
&=&\exp(x\sum_{i=0}^{i=m-1}2^{i})\\
\end{eqnarray*}
Note that by choosing a suitable sequence  $(E_k)$ and $(n_k)$, we
easily understand that   $$ \lim_{k\rightarrow \infty}\underset{x\in
E_{k}}{ess\,sup}\, \tilde{\omega}_{n_k}(x)=0.
$$ Ultimately,
$J_{T_{a,\omega}}(0)=L^p(\mathbb{R})$ by Theorem \ref{t3}. But
$T_{a,\omega}$ is not hypercyclic by \cite[Theorem 2.2]{chu}.
\end{exam}
\begin{exam}
Let $G=(\mathbb{R}, +)$, be the group of the real numbers equipped
with  Lebesgue measure.  Define the weight function $w$ on
$\mathbb{R}$ by
$$\omega(x)=\left\{
  \begin{array}{ll}
   2, & {x\leq -1}, \\
        -\frac{7}{4}x+\frac{1}{4}, & {-1< x < 0},\\
    \frac{7}{4}x+\frac{1}{4}, &{0<x\leq 1} ,\\
-\frac{7}{4}(x-1-m)+\frac{1}{4}, & {m< x \leq m+1 (m=1,3,5,...)},\\
\frac{7}{4}(x-m)+\frac{1}{4}, & {m< x \leq m+1 (m=2,4,6,...)}.
  \end{array}
\right.$$ Indeed, for each $x>2$ we adopt that
$\omega(x)=\omega(x-2)$. See the following figure:

%\includegraphics[width=0.5\textwidth]{a}

\begin{figure}[h]
  \centering
 \includegraphics[width=0.5\textwidth]{a}
  \caption{Weight function}
  \label{fig:10}
\end{figure}

 Fix a  nonzero number $a\in
\mathbb{R}$ and consider the convolution
$(\omega*\delta_{a})(x)=\omega(x- a)$, for each $x\in \mathbb{R}$
and $f\in L^p(\mathbb{R})$.   Fix $a=2$ and choose  compact subset
$\Delta\subseteq \mathbb{R}$ arbitrarily. Then for each $x\in
\Delta$ we have
$${\tilde{\omega}_{n}}(x)=\frac{1}{\prod_{i=0}^{i=n-1}\omega(xa^{-i})}=
\frac{1}{\prod_{i=0}^{i=n-1}\omega(x-2{i})} \rightarrow 0,$$ as
$n\rightarrow \infty$. Also, for each $x\in [0, \frac{1}{4}]$ we
have $\frac{1}{4}<\omega(x)<\frac{11}{16}$. Hence,
$$\omega_{n}(x)=\prod_{i=1}^{i=n}\omega(xa^{i})=
\prod_{i=1}^{i=n}\omega(x+2{i}) \rightarrow 0,$$ as $n\rightarrow
\infty$.



 Eventually, the conditions of Theorem \ref{t4} are satisfied  and  then
 guarantee that the corresponding weighted translation $T_{a,w}$ is $J$-class
on $L^p(\mathbb{R})$ but it is not hypercyclic by \cite[Theorem
2.2]{chu}.
 \end{exam}
\section{conclusions}
There are $J$-class weighted translations on locally compact groups
which are not hypercyclic. Indeed, according to Theorem \ref{t4}, we
may find many nonzero functions of $C_c(G)$ which are $J$-vectors
for non hypercyclic weighted translations on $L^p(G)$(\cite[Theorem
2.2]{chu}). Also,  $J$-class weighted translations  can occur on
non-separable $L^p(G)$ spaces whereas hypercyclic phenomena cannot.
As we have seen in the results, the extended limit sets i.e.,
$J_T(x)$ are playing powerful role rather than the orbits in the
dynamics of linear operators. For this, note that  the orbit of
weighted translation $T_{a,\omega}$ at zero cannot be whole space
$L^p(G)$ while it is possible $J_{T_{a,\omega}}(0)=L^p(G)$. Another
difference between $J$-sets and orbits of weighted translations on
locally compact groups, is in the case that  a right translated
element is torsion. In this case, unlike the case of non-dense
orbits of weighted translation operators (\cite[Lemma 1.1]{chu}), we
have $J_{T_{a,\omega}}(0)=L^p(G)$.







%-------------------------------------------------------------------------
\begin{thebibliography}{99}

\bibitem{ansari}
S. I. Ansari, Hypercyclic and cyclic vectors, J. Funct. Anal. 128 (1995), no. 2, 374-383.
%%%%%%%%%%%%%%%%

\bibitem{ak}
I. Akbarbaglu and M.R. Azimi, Universal family of translations on
weighted Orlicz spaces, Positivity  26 (2022), no. 1, Paper No. 1,
21 pp.
%%%%%%%%%%%%%%%%%%%%%%%%%%%%%%%%%%%%%%%%%%%%
\bibitem{kum} I. Akbarbaglu, M.R. Azimi, V. Kumar, Topologically transitive sequence
of cosine operators on Orlicz
spaces, Ann. Funct. Anal. 12 (2021), no. 1, Paper No. 3, 14 pp.

\bibitem{as}
M. Asadipour, B. Yousefi, On some properties of $J$-class operators.
Commun. Korean Math. Soc. 34 (2019), no. 1, 145-154.

%%%%%%%%%%%%%%%%
\bibitem{az3} M. R. Azimi,  $J$-class sequences of linear operators,
 Complex Anal. Oper. Theory 12(1) (2018), 293-303.

%\bibitem{az2}M. R. Azimi and I. Akbarbaglu, Hypercyclicity of weighted
%translations on Orlicz spaces, J. Oper. Matrices. 12 (2018), no. 1,27-37.
%%%%%%%%%%%%%%%%%%%%%%%%%%%%%%%%%%%%%%%
\bibitem{az} M. R. Azimi, M. Farmani, Subspace-hypercyclicity of conditional
 weighted translations
 on locally compact groups, Positivity 26, 58 (2022).

 \bibitem{azadv} M. R. Azimi, M. Farmani, Subspace-supercyclicity of conditional
 weighted type translations on $L^p(G)$,
 Adv. Oper. Theory 8, 37 (2023).

 \bibitem{azmu} M. R. Azimi,  V. M\"{u}ller,  A note on $J$-sets of linear
 operators,
 Rev. R. Acad. Cienc.
 Exactas F\'{\i}s. Nat. Ser. A Mat. RACSAM 105 (2011), no. 2, 449-453.
%%%%%%%%%%%%%%%%%%%%%%%%%%%%%%%%%%%%%%%%%%%%%


%%%%%%%%%%%%%%%%%%%%%%%%%%%%%%%%%%%%%%%%
\bibitem{baya}
F. Bayart and  \'{E}. Matheron, Dynamics of linear operators, Cambridge Tracts in Mathematics,
Cambridge University Press, Cambridge, (2009).

%%%%%%%%%%%%%%%%%%%%%%%%%%%%%%%%%%%%%%%%%%%%%%%%%%%%%%%%%%%%%5

\bibitem{chu}
C. Chen and C.-H. Chu, Hypercyclic weighted translations on groups, Proc. Amer. Math.
Soc. 139 (2011), no. 8, 2839-2846.
%%%%%%%%%%%%%%%%%%%%%%%%%%%%%%%%%%%%%%%%%%%%%%

\bibitem{chen}
C. C. Chen, Supercyclic and Ces$\grave{a}$ro hypercyclic weighted
translations on groups,
 Taiwanese J. Math. 16 (2012), no. 5, 1815-1827.
 %%%%%%%%%%%%%%%%%%%%%%%%%%%%%%%%%%%%%%%%
\bibitem{man}
  G. Costakis, A. Manoussos,  $J$-class operators and hypercyclicity,
   J. Operator Theory 67 (2012), no. 1, 101-119.

 \bibitem{cos}
G. Costakis, A. Manoussos, $J$-class weighted shifts on the space of
bounded sequences of complex numbers, Integral Equations Operator
Theory 62 (2008), no. 2, 149-158.


%%%%%%%%%%%%%%%%%%%%%%%%%%%%%%%%%%%%%%%%%%%%%%
\bibitem{gros}
K.-G. Grosse-Erdmann,  A. P. Manguillot, Linear chaos, Universitext.
Springer. London. (2011).




 %%%%%%%%%%%%%%%%%%%%%%%%%%%%%%%%%%


\bibitem{mar} H. Marzougui,  $J$-class abelian semigroups of matrices on $\Bbb{R}^{n}$,
 Appl. Math. Nonlinear Sci. 2(2) (2017),  519-528.


\bibitem{nas2} A. B. Nasseri, On the existence of $J$-class operators on Banach spaces,
 Proc. Amer. Math. Soc. 140 (2012), no. 10, 3549-3555.



 %%%%%%%%%%%%%%%%%%%%%%%%%%%%%%%%%%%%%%%%%
 \bibitem{sal}
 H. Salas, Hypercyclic weighted shifts,
 Trans. Amer. Math. Soc. 347 (1995), no. 3, 993-1004.

 %%%%%%%%%%%%%%%%%%%%%%%%%%%%%%%%%%%%%%






\end{thebibliography}
\end{document}
% ------------------------------------------------------------------------
